\section{Módulo de Punto de Venta (POS)}

El Punto de Venta (POS) está optimizado para agilizar la atención en el mostrador, reduciendo el tiempo de despacho a menos de 5 minutos y garantizando la coherencia entre las ventas y el descargo del almacén.

\subsection{Interfaz de Caja}

Para iniciar el flujo de venta en tienda, abra la sesión de su respectiva caja registradora. La pantalla principal del POS se divide en dos áreas principales: a la izquierda, el listado (ticket virtual) de los artículos que el cliente desea llevar; y a la derecha, el teclado numérico de cobro junto con las categorías de los artículos.

\vspace{1cm}
\begin{center}
\textit{[Espacio reservado para la captura de la Interfaz del POS]}
\end{center}
\vspace{1cm}

Pasos para registrar una venta:
\begin{enumerate}
    \item Seleccione los productos haciendo clic en ellos o utilizando el buscador (lupa).
    \item (Opcional) Modifique la cantidad, aplique algún descuento autorizado o seleccione un cliente específico de la base de datos si requiere comprobante con nombre.
    \item Haga clic en el botón \textbf{Pagos}.
    \item Seleccione el método de pago (Efectivo, Tarjeta, Transferencia) y valide el monto.
    \item Confirme la transacción. El sistema descontará inmediatamente los productos del stock.
\end{enumerate}

\subsection{Generación del Comprobante (Recibo)}

Al finalizar el pago, el sistema emitirá el recibo correspondiente, el cual puede ser impreso o enviado por correo electrónico.

\vspace{1cm}
\begin{center}
\textit{[Espacio reservado para la captura del Ticket/Recibo de Venta]}
\end{center}
\vspace{1cm}
